\scanpage{74}
\mspage{23}

\uline{Définition 1} Sous les conditions \unsure{équ.|équiv.} a) b) de la prop 1, on dit que $f$ est \uline{$W$-asphérique}. On dit $f$ est \uline{$W$-coasphérique} si $f^\circ$ est $W$-asphérique, i.e. si les $ {}_{y}\backslash X$ sont $W$-asphériques (pour $y\in\mathrm{Ob}\,Y$).

\deleted{La condition \unsure{a)}} \inserted{\uline{Dém}} Il suffit de prouver $\mathrm b\Rightarrow\mathrm a$. Par définition d'un isom local, on voit que la condition b) est stable par changement de base $Y'\longrightarrow Y$, avec $Y'$ un isom local (puisque les $Y'_{/y'}\xrightarrow{\sim}Y_{/y}$, donc $X'_{/y'}\simeq X_{/y}$). Il suffit donc de voir que si $f$ \unsure{satisfait} à b), alors $f\in W$. Mais c'est l'axiome $W(2\,\mathrm{bis})$.

\deleted{\diagram{\begin{tikzcd}[ampersand replacement=\&,column sep=tiny,row sep=tiny] \&\&\deleted{X}\arrow[d]\\ \deleted{\unsure{Y_{/y'}=Y'}}\arrow[r]\&\deleted{\unsure{Y}}\&\deleted{Y}\&\deleted{Y'}\arrow[l]\\ \deleted{Y'_{/y'}}\arrow[u]\arrow[r]\&\deleted{Y_{/y}}\arrow[u]\&\deleted{Y_{/y}}\arrow[u]\&\deleted{Y'_{/y'}}\arrow[l]\arrow[u] \end{tikzcd}}} \diagram{\begin{tikzcd}[ampersand replacement=\&,column sep=small,row sep=small] X\arrow[d]\&\&\\ Y\&Y'\arrow[l]\&\\ \&Y_{/y}\arrow[ul]\&Y'_{/y'}\arrow[ul]\arrow[l] \end{tikzcd}}

\uline{Remarque} Si $f$ a un adjoint à droite $g$, alors pour tout $y\in Y$, $X_{/y}\simeq X_{/g(y)}$, donc $X_{/y}$ a un objet final et est donc $W$-asphérique. Donc si $f$ a un adjoint à droite (à gauche), alors $f$ est $W$-asphérique ($W$-coasphérique).

Mais $=$ si $f$ n'a pas d'adjoint à droite, il y a \unsure{cepdt} $=$ un adjoint généralisé \[g:Y\longrightarrow X^\wedge,\] et on a \[X_{/y}\simeq X_{/g(y)}\qquad(\text{pour l'inclusion tt. fid. }X\deleted{\unsure{{}^{\wedge}}}\hookrightarrow X^\wedge),\] donc $\mathrm{Top}(X_{/y})$ s'identifie au topos induit $\mathrm{Top}(X)_{/g(y)}$. La condition $f$ $W$-asphérique signifie donc que les \inserted{faisceaux} $g(y)$ \deleted{dans} $\mathrm{Top}(X)$ sont

% archive sheet number: 74
